\documentclass[border=10pt]{standalone}
\usepackage[no-math]{fontspec}
\usepackage[brazil]{babel}
\usepackage{amsmath,amssymb}
\IfFileExists{C:/Windows/Fonts/Lato_Regular.ttf}{%
  \setmainfont{Lato}[
    Path = C:/Windows/Fonts/,
    UprightFont = Lato_Regular.ttf,
    BoldFont = Lato_Bold.ttf,
    ItalicFont = Lato_Italic.ttf,
    BoldItalicFont = Lato_BoldItalic.ttf
  ]
  \setsansfont{Lato}[
    Path = C:/Windows/Fonts/,
    UprightFont = Lato_Regular.ttf,
    BoldFont = Lato_Bold.ttf,
    ItalicFont = Lato_Italic.ttf,
    BoldItalicFont = Lato_BoldItalic.ttf
  ]
}{%
  \RequirePackage[default]{lato}
}
\renewcommand{\familydefault}{\sfdefault}

\usepackage{xcolor}
\usepackage{tikz}
\usepackage{pgfplots}
\pgfplotsset{compat=1.18}
\usetikzlibrary{arrows.meta,positioning,shapes.geometric,fit,backgrounds,calc,decorations.pathreplacing,patterns}

% Paleta institucional PLANTERR / UEFS (alinhada ao estilo planterr-slides.css / 09a)
\definecolor{etNavy}{HTML}{1F4268}
\definecolor{etNavyDark}{HTML}{152E4A}
\definecolor{etGold}{HTML}{B38A3D}
\definecolor{etGoldDark}{HTML}{8C6927}
\definecolor{etTeal}{HTML}{0E7C7B}
\definecolor{etForest}{HTML}{2D6A4F}
\definecolor{etAlert}{HTML}{9E2A2B}
\definecolor{etInk}{HTML}{26313D}
\definecolor{etMuted}{HTML}{5B6573}
\definecolor{etBorder}{HTML}{D5DEE8}

\definecolor{bgNavy}{HTML}{EDF2F7}
\definecolor{bgGold}{HTML}{F8F3E9}
\definecolor{bgTeal}{HTML}{EBF6F5}
\definecolor{bgForest}{HTML}{EEF5F0}
\definecolor{bgAlert}{HTML}{FDF2F2}
\definecolor{bgCanvas}{HTML}{FFFFFF}

\tikzset{
  >=Stealth,
  etBox/.style={
    draw=#1, fill=white, rounded corners=5pt, line width=1.15pt,
    inner sep=7pt, align=center, font=\footnotesize\sffamily, text=etInk
  },
  etTint/.style={
    draw=#1, fill=#1!8!white, rounded corners=5pt, line width=1.15pt,
    inner sep=7pt, align=center, font=\footnotesize\sffamily, text=etInk
  },
  etBadge/.style={
    fill=#1, text=white, rounded corners=9pt, inner xsep=8pt, inner ysep=3pt,
    font=\scriptsize\bfseries\sffamily
  },
  etTitle/.style={font=\normalsize\bfseries\sffamily, text=etNavy, align=center},
  etSub/.style={font=\footnotesize\sffamily, text=etMuted, align=center},
  etArrow/.style={->, >=Stealth, line width=1.35pt},
  etDArrow/.style={<->, >=Stealth, line width=1.35pt}
}


\begin{document}
\begin{tikzpicture}[font=\sffamily]

% Título
\node[etTitle, font=\large\bfseries\sffamily, text=etNavy] (title) at (7.75, 6.85)
  {Rubrica Oficial AV3 (50\%) e Tripé Integrador do Projeto Final};

% Zona esquerda: Painel Donut da Rubrica AV3
\node[
  draw=etBorder, fill=white, rounded corners=8pt, line width=1.15pt,
  minimum width=6.60cm, minimum height=6.40cm, anchor=north west
] (leftPanel) at (0.30, 6.20) {};

\node[font=\small\bfseries\sffamily, text=etNavy, anchor=north] at ([yshift=-0.26cm]leftPanel.north)
  {Composição da Avaliação Final (AV3)};

% Donut ring centrado em (3.60, 3.85) com folga superior e inferior
% 30% = 108 deg, 30% = 108 deg, 25% = 90 deg, 15% = 54 deg
\begin{scope}[shift={(3.60, 3.85)}]
  \def\Rin{0.76}
  \def\Rout{1.44}
  \def\Rmid{1.10}

  % Fatia 1: Leitura Territorial 30% (90 -> -18)
  \fill[etNavy] (90:\Rout) arc (90:-18:\Rout) -- (-18:\Rin) arc (-18:90:\Rin) -- cycle;
  \draw[white, line width=1.6pt] (90:\Rin) -- (90:\Rout);
  \node[font=\footnotesize\bfseries\sffamily, text=white] at (36:\Rmid) {30\%};

  % Fatia 2: Estratégia & Modelo 30% (-18 -> -126)
  \fill[etTeal] (-18:\Rout) arc (-18:-126:\Rout) -- (-126:\Rin) arc (-126:-18:\Rin) -- cycle;
  \draw[white, line width=1.6pt] (-18:\Rin) -- (-18:\Rout);
  \node[font=\footnotesize\bfseries\sffamily, text=white] at (-72:\Rmid) {30\%};

  % Fatia 3: Viabilidade & Valuation 25% (-126 -> -216)
  \fill[etGold] (-126:\Rout) arc (-126:-216:\Rout) -- (-216:\Rin) arc (-216:-126:\Rin) -- cycle;
  \draw[white, line width=1.6pt] (-126:\Rin) -- (-126:\Rout);
  \node[font=\footnotesize\bfseries\sffamily, text=white] at (-171:\Rmid) {25\%};

  % Fatia 4: Comunicação & Defesa 15% (-216 -> -270)
  \fill[etAlert] (-216:\Rout) arc (-216:-270:\Rout) -- (-270:\Rin) arc (-270:-216:\Rin) -- cycle;
  \draw[white, line width=1.6pt] (-216:\Rin) -- (-216:\Rout);
  \node[font=\footnotesize\bfseries\sffamily, text=white] at (-243:\Rmid) {15\%};

  % Centro do donut
  \node[align=center, font=\small\bfseries\sffamily, text=etNavy] at (0, 0.11) {AV3};
  \node[align=center, font=\scriptsize\sffamily, text=etMuted] at (0, -0.22) {50\% da nota};
\end{scope}

% Legenda em barras limpas abaixo do donut (y = 1.85 até 0.29, painel vai até -0.20)
\tikzset{
  legRow/.style={
    draw=#1!40, fill=#1!7!white, rounded corners=4pt, line width=0.8pt,
    minimum width=5.75cm, minimum height=0.43cm, inner sep=0pt
  }
}
\node[legRow=etNavy]  (l1) at (3.60, 1.85) {};
\fill[etNavy, rounded corners=2pt] ([xshift=0.14cm, yshift=-0.10cm]l1.north west) rectangle ([xshift=0.38cm, yshift=0.10cm]l1.south west);
\node[anchor=west, font=\scriptsize\bfseries\sffamily, text=etInk] at ([xshift=0.50cm]l1.west) {Leitura Territorial};
\node[anchor=east, font=\scriptsize\bfseries\sffamily, text=etNavy] at ([xshift=-0.20cm]l1.east) {30\%};

\node[legRow=etTeal]  (l2) at (3.60, 1.33) {};
\fill[etTeal, rounded corners=2pt] ([xshift=0.14cm, yshift=-0.10cm]l2.north west) rectangle ([xshift=0.38cm, yshift=0.10cm]l2.south west);
\node[anchor=west, font=\scriptsize\bfseries\sffamily, text=etInk] at ([xshift=0.50cm]l2.west) {Estratégia \& Modelo};
\node[anchor=east, font=\scriptsize\bfseries\sffamily, text=etTeal] at ([xshift=-0.20cm]l2.east) {30\%};

\node[legRow=etGold]  (l3) at (3.60, 0.81) {};
\fill[etGold, rounded corners=2pt] ([xshift=0.14cm, yshift=-0.10cm]l3.north west) rectangle ([xshift=0.38cm, yshift=0.10cm]l3.south west);
\node[anchor=west, font=\scriptsize\bfseries\sffamily, text=etInk] at ([xshift=0.50cm]l3.west) {Viabilidade \& Valuation};
\node[anchor=east, font=\scriptsize\bfseries\sffamily, text=etGoldDark] at ([xshift=-0.20cm]l3.east) {25\%};

\node[legRow=etAlert] (l4) at (3.60, 0.29) {};
\fill[etAlert, rounded corners=2pt] ([xshift=0.14cm, yshift=-0.10cm]l4.north west) rectangle ([xshift=0.38cm, yshift=0.10cm]l4.south west);
\node[anchor=west, font=\scriptsize\bfseries\sffamily, text=etInk] at ([xshift=0.50cm]l4.west) {Comunicação \& Defesa};
\node[anchor=east, font=\scriptsize\bfseries\sffamily, text=etAlert] at ([xshift=-0.20cm]l4.east) {15\%};

% Zona direita: Arquitetura vertical em 3 etapas integradoras
\node[
  draw=etBorder, fill=white, rounded corners=8pt, line width=1.15pt,
  minimum width=7.95cm, minimum height=6.40cm, anchor=north west
] (rightPanel) at (7.25, 6.20) {};

\node[font=\small\bfseries\sffamily, text=etNavy, anchor=north] at ([yshift=-0.26cm]rightPanel.north)
  {Tripé Integrador do Projeto Final};

\tikzset{
  triCard/.style={
    draw=#1, fill=#1!6!white, rounded corners=6pt, line width=1.2pt,
    minimum width=7.15cm, minimum height=1.22cm, inner xsep=11pt, inner ysep=8pt,
    align=left, font=\footnotesize\sffamily, text=etInk, text width=6.60cm
  }
}

\node[triCard=etNavy] (u1) at (11.225, 4.62) {\raggedright\hyphenpenalty=10000%
  \textcolor{etNavy}{\textbf{Unidade I \textperiodcentered~Diagnóstico Territorial}}\\[4pt]
  \textit{Superar isomorfismo \textperiodcentered~ler APL, SRI e Hélice \mbox{Quíntupla}}
};

\node[triCard=etTeal] (u2) at (11.225, 2.75) {\raggedright\hyphenpenalty=10000%
  \textcolor{etTeal}{\textbf{Unidade II \textperiodcentered~Estratégia \& Modelo}}\\[4pt]
  \textit{Oportunidade, PI, Pitch Canvas e Fomento Regional}
};

\node[triCard=etGold] (u3) at (11.225, 0.88) {\raggedright\hyphenpenalty=10000%
  \textcolor{etGoldDark}{\textbf{Unidade III \textperiodcentered~Viabilidade \& Retenção}}\\[4pt]
  \textit{Localização, VPL/TIR, Valuation e IRV local}
};

% Setas verticais conectoras
\draw[etArrow, etNavy] (u1.south) -- (u2.north);
\draw[etArrow, etTeal] (u2.south) -- (u3.north);

\end{tikzpicture}
\end{document}
